\documentclass{beamer} % DocumentClass = Presentation.
\usetheme{Boadilla} % Theme of the presentation.
%\usecolortheme{dolphin} % Theme color palette.
%\useinnertheme{metropolis} % Style of inner elements.
%\useoutertheme{miniframes} % Style of outer elements.
%\usefonttheme{serif} % Some font styles.
% All combinations of themes and colors at
% https://hartwork.org/beamer-theme-matrix/

% FONTS OF THE DOCUMENT.
\usepackage{fontspec}
% Font corpo = Roboto Condensed. 
% Font titolo = Rotis Semi Serif.


% Presentation Details.
\title{SNS Beamer Template}
\author{Andrea Canzio}
\institute[]{\includegraphics[width=5cm]{images/logo_SNS_blue.png}}
\logo{\includegraphics[width=3cm]{images/logo_SNS_blue.png}\vspace{193pt}}
\date{\today}

% Packages.
\usepackage{graphicx}


% Begin of the document.
\begin{document}

\begingroup
\logo{}
\begin{frame}
    \titlepage
\end{frame}
\endgroup

% Recurrent Frame.
\AtBeginSection[]{
    \begin{frame}{Outline}
        \tableofcontents[currentsection, hideallsubsections]
    \end{frame}
}

\section{Introduction}
\subsection{The very beginning}

\begin{frame}{First Frame Title}
    blablabla...
\end{frame}

\begin{frame}
    \frametitle{Another way to set the title}
    blabla again.
\end{frame}

\section{Second one}

\begin{frame}{Yet another frame}
    Something interesting.
\end{frame}

\begin{frame}
    \frametitle{List}
    \begin{itemize}
        \item Point A
        \item Point B
        \begin{itemize}
            \item part 1
            \item part 2
        \end{itemize}
        \item Point C
        \item Point D
    \end{itemize}
    \begin{enumerate}[I]
        \item Point A
        \item Point B
        \begin{enumerate}[i]
            \item part 1
            \item part 2
        \end{enumerate}
        \item Point C
        \item Point D
    \end{enumerate}
\end{frame}

\begin{frame}
    \frametitle{Using Columns}
    \begin{columns}
        \column{0.5\textwidth}
            some text
        \column{0.5\textwidth}
            some other text
    \end{columns}
\end{frame}

\begin{frame}
    \frametitle{Pictures}
    \begin{figure}
        \includegraphics[scale=0.03]{images/image_placeholder.png}
        \caption{Picture of a Lama}
    \end{figure}
    some text.
\end{frame}

\begin{frame}
    \frametitle{Pictures part 2}
    \begin{columns}
        \column{0.5\textwidth}
            \begin{figure}
                \includegraphics[scale=0.03]{images/image_placeholder.png}
                \caption{Picture of a Lama}
            \end{figure}
        \column{0.5\textwidth}
        some text.
    \end{columns}
\end{frame}

\begin{frame}
    \frametitle{Description, useful for definitions}
    \begin{description}
        \item[API] Application Programming Interface
        \item[LAN] Local Area Network
        \item[ASCII] American Standard Code for Information Interchange
    \end{description}
\end{frame}

\begin{frame}
    \frametitle{Table}
    \begin{table}
        \begin{tabular}{l | c | c | c | c }
            Competitor Name & Swim & Cycle & Run & Total \\
            \hline \hline
            John T & 13:04 & 24:15 & 18:34 & 55:53 \\ 
            Norman P & 8:00 & 22:45 & 23:02 & 53:47\\
            Alex K & 14:00 & 28:00 & n/a & n/a\\
            Sarah H & 9:22 & 21:10 & 24:03 & 54:35 
        \end{tabular}
        \caption{Triathlon results}
    \end{table}
\end{frame}

\begin{frame}
    \frametitle{Frame with a Block}
    \begin{block}{Block Title}
        Lorem ipsum dolor sit amet, consectetur adipisicing elit, 
        sed do eiusmod tempor incididunt ut labore et 
        dolore magna aliqua.
    \end{block}
    \begin{alertblock}{Block Title}
        Lorem ipsum dolor sit amet, consectetur adipisicing elit, 
        sed do eiusmod tempor incididunt ut labore et 
        dolore magna aliqua.
    \end{alertblock}
    \begin{example}
        Lorem ipsum dolor sit amet, consectetur adipisicing elit, 
        sed do eiusmod tempor incididunt ut labore et
        dolore magna aliqua.
    \end{example}
\end{frame}

\begin{frame}
    \frametitle{Some other kinds of blocks}
    \begin{definition}
        A prime number is a number that...
    \end{definition}
    \begin{theorem}[Pythagoras] 
        $ a^2 + b^2 = c^2$
    \end{theorem}
    \begin{corollary}
        $ x + y = y + x  $
    \end{corollary}
    \begin{proof}
        $\omega +\phi = \epsilon $
    \end{proof}
\end{frame}

\begin{frame}[fragile]
    % The fragile option help with newlines and other special characters
    % that now are read explicitly. Tabs are also now read explicitly!
    \frametitle{Bad code example}
    \begin{semiverbatim}
\\begin\{frame\}
    \\frametitle\{Outline\}
    \\tableofcontents
        Other text to make you notice the explicit TABS.
\\end\{frame\}
    \end{semiverbatim}
\end{frame}

\begin{frame}
    \label{FrameToGo}
    \frametitle{This is the labelled frame}
    This is a random frame with a label.
\end{frame}

\begin{frame}
    \frametitle{Frame with the reference to the labelled one}
    Hyperlinks in different styles:\\
    \hyperlink{FrameToGo}{Hyperlink without button}\\
    \hyperlink{FrameToGo}{\beamerbutton{General style button}}\\
    \hyperlink{FrameToGo}{\beamergotobutton{Go to style button}}\\
    \hyperlink{FrameToGo}{\beamerskipbutton{Skip style button}}\\
    \hyperlink{FrameToGo}{\beamerreturnbutton{Return style button}}
\end{frame}

\begin{frame}
    \frametitle{Frame with a Paused list}
    \begin{itemize}
        \pause
        \item I really
        \pause
        \item don't like
        \pause
        \item to pause a frame.
        \pause
        \item You can be better than this.
    \end{itemize}
\end{frame}

\begin{frame}
    \frametitle{Overlays}
    Overlays are a bit better then \\pause, in the sense that make you able to do more complicated stuff like the following:\\
    \onslide<2->{First Line of Text}\\
    \onslide<4->{Second Line of Text}\\
    \onslide<3->{Third Line of Text}
\end{frame}

\begin{frame}
    \frametitle{Overlays with lists}
    Overlays also works with lists.
    Here we also show some more complicated way of using them.
    \begin{enumerate}[(I)]
        \item<1-> Point A
        \item<2-> Point B
        \begin{itemize}
            \item<3-> part 1
            \item<4-> part 2
        \end{itemize}
        \item<2,4,6-> Appears on slide 2, 4, $\geq$6.
        \item<5-> Point D
    \end{enumerate}
\end{frame}

\begin{frame}
    \frametitle{Overlays}
    \only<1>{First Line of Text}
    \only<2>{Second Line of Text}
    \only<3>{Third Line of Text}
\end{frame}

\begin{frame}
    \frametitle{Other Overlays effects}
    \textbf<2>{Example Text}
    \textit<2>{Example Text}
    \textsl<2>{Example Text}
    \textrm<2>{Example Text}
    \textsf<2>{Example Text}
    \textcolor<2>{orange}{Example Text}
    \alert<2>{Example Text}
    \structure<2>{Example Text}
\end{frame}

\begin{frame}
    \frametitle{Maths Blocks}
    \begin{theorem}<1->[Pythagoras] 
        $ a^2 + b^2 = c^2$
    \end{theorem}
    \begin{corollary}<3->
        $ x + y = y + x  $
    \end{corollary}
    \begin{proof}<2->
        $\omega +\phi = \epsilon $
    \end{proof}
\end{frame}


\setbeamercovered{transparent} %If you want the covered text to be transparent instead of invisibile!
\begin{frame}
\frametitle{Tables}
\begin{table}
\begin{tabular}{l | c | c | c | c }
Competitor Name & Swim & Cycle & Run & Total \\
\hline \hline
John T & 13:04 & 24:15 & 18:34 & 55:53 \onslide<2-> \\ 
Norman P & 8:00 & 22:45 & 23:02 & 53:47 \onslide<3->\\
Alex K & 14:00 & 28:00 & n/a & n/a \onslide<4->\\
Sarah H & 9:22 & 21:10 & 24:03 & 54:35 
\end{tabular}
\caption{Triathlon results}
\end{table}
\end{frame}




\end{document}
